\documentclass[12pt]{article}
\usepackage[utf8]{inputenc}
\usepackage[T1]{fontenc}
\usepackage{hyperref}
\title{Huarij\'io Language Resources for K-12 Curriculum Development}
\author{Compiled by Claude Sonnet 5 (claude-sonnet-5) for Andr\'as Kornai}
\date{}
\begin{document}
\maketitle

\section{Introduction}
Huarij\'io (Guarij\'io; ISO 639-3 \texttt{var}, Glottocode huar1255) is a Uto-Aztecan language of Sonora and Chihuahua, Mexico. Two varieties are described: Mountain Guarij\'io and River Guarij\'io. Searches found no school curriculum materials; the resources below are reference works only.

\section{Primary Educational Resources}
No primers, textbooks or school programmes were found in the search results. Omniglot mentions Huarij\'io-language radio programming from a station in Etchojoa, Sonora, but gives no teaching materials.

\section{Linguistic and Reference Materials}
\begin{itemize}
\item Wick R. Miller (1996), \emph{Guarij\'io: Gram\'atica, Textos y Vocabulario}, a grammar with texts and vocabulary in Spanish (catalogue entry via Glottolog): \url{https://glottolog.org/resource/reference/id/665425}
\item Grambank entry for Huarij\'io, which cites Miller (1996) and a grammar of River Warihio by Armend\'ariz (2005): \url{https://grambank.clld.org/languages/huar1255}
\item Glottolog language page: \url{https://glottolog.org/resource/languoid/id/huar1255}
\item Omniglot page with alphabet chart and sample text: \url{https://omniglot.com/writing/huarijio.htm}
\item Endangered Languages Project bibliography: \url{https://www.endangeredlanguages.com/lang/6203/bibliography?hl=es}
\item Wikipedia overview: \url{https://en.wikipedia.org/wiki/Huarijio_language}
\end{itemize}

\section{Religious and Community Resources}
None were found in the search results. A general page on the Guarij\'io people exists: \url{https://en.wikipedia.org/wiki/Guarijio_people}

\section{Digital and Audio Resources}
No audio, video or software resources were identified in the search results beyond the radio mention at Omniglot (\url{https://omniglot.com/writing/huarijio.htm}).

\section{Curriculum Development Context}
The language is reported as not taught in schools. The only substantial resources found are Miller's grammar, texts and vocabulary (Spanish-language) and Armend\'ariz's grammar of the River variety. A curriculum would need to start from these, from the community, and from new materials. Note that resources may cover only one of the two varieties.
\end{document}